\chapter{Project Plan}

% ──────────────────────────────────────────
%  HEADER
% ──────────────────────────────────────────
\begin{center}
    {\Large\textbf{Project Plan}}\\[0.4em]
    \textit{Project: GG-Code -- Domain-Specific Language for Generating and Editing G-Code}\\[0.2em]
    \textit{Institution: Technical University of Moldova}\\[0.2em]
    \textit{Faculty: Computers, Informatics and Microelectronics}\\[0.2em]
    \textit{Department: Software Engineering and Automatics}\\[0.2em]
    \textit{Group: FAF-241 \textbar\ PBL Team No. 4 \textbar\ 2026}\\[0.2em]
    \textit{Mentor: lect. univ. Irina Cojuhari}
\end{center}

\vspace{1em}

% ──────────────────────────────────────────
%  1. PROJECT OVERVIEW
% ──────────────────────────────────────────
\section*{1. Project Overview}

GG-Code is a domain-specific language developed as part of the ELSD course at the Technical
University of Moldova. The project was initiated by PBL Team No.~4, group FAF-241, and the
domain was officially chosen at the first team meeting on 29 January 2026.

The goal of the project is to create a language that sits between existing CAD/CAM software
and the low-level G-code instructions used to control CNC machines, 3D printers, and laser
engravers. Instead of working directly with cryptic numeric machine instructions, users can
write clear, human-readable code that the compiler then translates into standard G-code.

% ──────────────────────────────────────────
%  2. PROBLEM
% ──────────────────────────────────────────
\section*{2. Problem}

G-code has been the industry standard for controlling manufacturing machines since the 1950s.
While it works, it was never designed to be written or edited by hand. A single print job can
produce tens of thousands of lines of code that are nearly impossible to read or safely modify
without specialist knowledge.

At the same time, the CAD/CAM tools that generate this code treat it as a black box. If a user
needs to make even a small adjustment, they have no easy way to do it -- they either have to
regenerate everything from scratch or risk editing raw machine instructions by hand. There is
no middle ground that gives users both control and clarity.

% ──────────────────────────────────────────
%  3. PROPOSED SOLUTION
% ──────────────────────────────────────────
\section*{3. Proposed Solution}

The project proposes a domain-specific language called GG-Code that fills the gap between
high-level design tools and low-level machine instructions. The language is designed to be
readable and writable by people, while still compiling down to valid G-code that any standard
machine can execute.

GG-Code supports variables, loops, conditions, and reusable operations, which means users can
describe what they want the machine to do in a natural way without having to think about
individual motor movements. The compiler handles the translation and checks for errors before
anything reaches the machine.

The language is formally defined using a context-free grammar and implemented with ANTLR4.
The project includes a complete lexer, parser, and abstract syntax tree, along with a test
suite and a full project report written in \LaTeX.

% ──────────────────────────────────────────
%  4. TEAM
% ──────────────────────────────────────────
\section*{4. Team}

The project is developed by four students from group FAF-241 under the supervision of mentor
Irina Cojuhari. Roles and tasks were distributed at the first team meeting on 29 January 2026.

\begin{itemize}[leftmargin=2em, label={--}]
    \item Andrei Cliucovschi
    \item Sofia Zhuchikova
    \item Alexei Luchiciov
    \item Dragos Nedelea
\end{itemize}

% ──────────────────────────────────────────
%  FOOTER
% ──────────────────────────────────────────
\vspace{2em}
\begin{center}
    \small\textit{GG-Code Project Plan \textbar\ PBL Team 4 \textbar\ FAF-241 \textbar\ Technical University of Moldova \textbar\ 2026}
\end{center}

